\chapter{滕文公章句下}
\kaishi*

%滕文公章句下

\section{陳代勸見諸侯}

陳代曰：「不見諸侯，宜若小然。今一見之，大則以王，小則以霸。且志曰：『枉尺而直尋。』宜若可爲也。」\jiazhu[1]{陳代，孟子弟子也。代見諸侯有來聘請孟子，孟子有所不見，以爲孟子欲以是爲介，故言此介得無爲狹小乎？如一見之，儻得行道，可以輔致霸王乎？志，記也。枉尺直尋，欲使孟子屈己信道，故言宜若可爲也。}

孟子曰：「昔齊景公田，招虞人以旌，不至，將殺之。」\jiazhu[1]{虞人，守苑囿之吏也。招之當以皮冠，而以旌，故不至也。}

「志士不忘在溝壑，勇士不忘喪其元。孔子奚取焉？取非其招不往也。如不待其招而往，何哉？」\jiazhu[1]{志士，守義者也。君子固窮，故常念死無棺槨，沒溝壑而不恨也。勇士，義勇者也。元，首也。以義則喪首不顧也。孔子奚取？取守死善道，非禮招己則不往。言虞人不得其招尚不往，如何君子而不待其招，直事妄見諸侯者，何爲也？}

「且夫枉尺而直尋者，以利言也。如以利，則枉尋直尺而利，亦可爲與？」\jiazhu[1]{尺小尋大，不可枉大就小，而要以利也。}

「昔者趙簡子使王良與嬖奚乘，終日而不獲一禽。嬖奚反命曰：『天下之賤工也。』」\jiazhu[1]{趙簡子，晉卿也。王良，善御者也。嬖奚，簡子幸臣。以不能得一禽，故反命於簡子，謂王良天下鄙賤之工師也。}

「或以告王良。良曰：『請復之。』」\jiazhu[1]{聞嬖奚賤之，故請復與乘。}

「強而後可。」\jiazhu[1]{強嬖奚乃肯行。}

「一朝而獲十禽。嬖奚反命曰：『天下之良工也。』」\jiazhu[1]{以一朝得十禽，故謂之良工。}

「簡子曰：『我使掌與女乘。』謂王良。良不可。」\jiazhu[1]{掌，主也。使王良主與女乘。}

「曰：『吾爲之範我馳驅，終日不獲一；爲之詭遇，一朝而獲十。』」\jiazhu[1]{範，法也。王良曰：我爲之法度之御，應禮之射，正殺之禽，不能得一。橫而射之曰詭遇。非禮之射，則能獲十。言嬖奚小人，也不習於禮。}

「《詩》云：『不失其馳，舍矢如破。』我不貫與小人乘，請辭。」\jiazhu[1]{詩小雅車攻之篇也。言御者不失其馳驅之法，則射者必中之。順毛而入，順毛而出，一發貫臧，應矢而死者如破矣。此君子之射也。貫，習也。我不習與小人乘，不願掌與嬖奚同乘，故請辭。}

「御者且羞與射者比，比而得禽獸，雖若丘陵，弗爲也。如枉道而從彼，何也？」\jiazhu[1]{孟子引此以喻陳代云：御者尚知羞恥此射者，不欲與比，子如何欲使我枉正道而從彼驕慢諸侯而見之？}

「且子過矣！枉己者，未有能直人者也。」\jiazhu[1]{謂陳代之言過謬也。人當以直矯枉耳，己自枉曲，何能正人？}\par \jiazhu[1]{章指言：脩禮守正，非招不往。枉道富貴，君子不許。是以諸侯雖有善其辭命，伯夷亦不屑就也。}

\section{景春論大丈夫}

景春曰：「公孫衍、張儀豈不誠大丈夫哉？一怒而諸侯懼，安居而天下熄。」\jiazhu[1]{景春，孟子時人，爲從橫之術者。公孫衍，魏人也，號爲犀首，常佩五國相印，爲從長，秦王之孫，故曰公孫。張儀，合從者也。一怒則構諸侯，使彊陵弱，故言懼也。安居不用辭說，則天下兵革熄也。}

孟子曰：「是焉得爲大丈夫乎？子未學禮乎？丈夫之冠也，父命之；女子之嫁也，母命之，往送之門，戒之曰：『往之女家，必敬必戒，無違夫子！』以順爲正者，妾婦之道也。」\jiazhu[1]{孟子以禮言之，男子之道當以義匡君，女子則當婉順從人耳。男子之冠則命曰就爾成德。今此二子從君順指，行權合從，無輔弼之義，安得爲大丈夫也？}

「居天下之廣居，立天下之正位，行天下之大道。得志，與民由之；不得志，獨行其道。富貴不能淫，貧賤不能移，威武不能屈。此之謂大丈夫。」\jiazhu[1]{廣居，謂天下也。正位，謂男子純乾正陽之位也。大道，仁義之道也。得志行正與民共之，不得志隱居獨善其身，守道不回也。淫，亂其心也。移，易其行也。屈，挫其志也。三者不惑，乃可謂大丈夫。}\par \jiazhu[1]{章指言：以道匡君，非禮不運，稱大丈夫。阿意用謀，善戰務勝，事雖有剛，心歸柔順，故云妾婦，以況儀、衍。}

\section{周霄問君子仕否}

周霄問曰：「古之君子仕乎？」\jiazhu[1]{周霄，魏人。問君子之道當仕否。}

孟子曰：「仕。《傳》曰：『孔子三月無君，則皇皇如也，出疆必載質。』」\jiazhu[1]{質，臣所執以見君者也。三月，一時也。物變而不佐君化，故皇皇如有所求而不得。}

「公明儀曰：『古之人三月無君，則弔。』」\jiazhu[1]{公明儀，賢者也。而言古人三月無君則弔，明當仕也。}

「三月無君則弔，不以急乎？」\jiazhu[1]{周霄怪乃弔於三月無君，何其急也？}

曰：「士之失位也，猶諸侯之失國家也。《禮》曰：『諸侯耕助，以供粢盛；夫人蠶繅，以爲衣服。犧牲不成，粢盛不絜，衣服不備，不敢以祭。惟士無田，則亦不祭。牲殺、器皿、衣服不備，不敢以祭，則不敢以宴。』亦不足弔乎？」\jiazhu[1]{諸侯耕助者，躬耕勸率其民，收其藉助，以供粢盛。粢，稷；盛，稻也。夫人親織蠶繅之事，以率女功衣服祭服。不成，不實肥腯也。惟，辭也。言惟絀禄之士無圭田者不祭。牲必特殺，故曰殺。皿，所以覆器者也。不祭則不宴，猶喪人也，不亦可弔乎？}

「出疆必載質，何也？」\jiazhu[1]{周霄問出疆何爲復載質。}

曰：「士之仕也，猶農夫之耕也。農夫豈爲出疆舍其耒耜哉？」\jiazhu[1]{孟子言仕之爲急，若農夫不可不耕。}

曰：「晉國亦仕國也，未嘗聞仕如此其急。仕如此其急也，君子之難仕，何也？」\jiazhu[1]{魏本晉也，故周霄曰我晉人也。亦仕而不知其急若此。若此君子何爲難仕？君子謂孟子，何爲不急仕也？}

曰：「丈夫生而願爲之有室，女子生而願爲之有家。父母之心，人皆有之。不待父母之命、媒妁之言，鑽穴隙相窺，踰牆相從，則父母、國人皆賤之。」\jiazhu[1]{言人不可觸情從欲，須禮而行。}

「古之人未嘗不欲仕也，又惡不由其道。不由其道而往者，與鑽穴隙之類也。」\jiazhu[1]{言古之人雖欲仕，如不由其正道，是與鑚穴者何異？}\par \jiazhu[1]{章指言：君子務仕，思播其道，達義行仁，待禮而動。苟容干禄，踰牆之女，人之所賤，故弗爲也。}

\section{彭更問傳食諸侯}

彭更問曰：「後車數十乘，從者數百人，以傳食於諸侯，不以泰乎？」\jiazhu[1]{泰，甚也。彭更，孟子弟子，怪孟子徒衆多而傳食於諸侯之國，得無爲甚奢泰也。}

孟子曰：「非其道，則一簞食不可受於人；如其道，則舜受堯之天下，不以爲泰。子以爲泰乎？」\jiazhu[1]{簞，笥也。非以其道，一笥之食不可受也。子以舜受堯天下爲泰乎？}

曰：「否。士無事而食，不可也。」\jiazhu[1]{彭更曰：不以舜爲泰也。謂仕無功事而虛食人者，不可也。}

曰：「子不通功易事，以羨補不足，則農有餘粟，女有餘布。子如通之，則梓匠輪輿皆得食於子。」\jiazhu[1]{孟子言凡人當通功易事，乃可各以奉其用。梓、匠，木工也；輪人、輿人，作車者。交易則得食於子之所有矣。周禮攻木之工七，梓、匠、輪、輿是其四者。羨，餘也。}

「於此有人焉，入則孝，出則悌，守先王之道，以待後之學者，而不得食於子。子何尊梓匠輪輿而輕爲仁義者哉？」\jiazhu[1]{入則事親孝，出則敬長順也。悌，順也。守先王之道，上德之士，可以化俗者。若此不得食子之禄，子何尊彼而賤此也？}

曰：「梓匠輪輿，其志將以求食也。君子之爲道也，其志亦將以求食與？」\jiazhu[1]{彭更以爲彼志於食，此亦但志食也。}

曰：「子何以其志爲哉？其有功於子，可食而食之矣。且子食志乎？食功乎？」\jiazhu[1]{孟子言禄以食功，子何食乎？}

曰：「食志。」\jiazhu[1]{彭更以爲當食志也。}

曰：「有人於此，毀瓦畫墁，其志將以求食也，則子食之乎？」\jiazhu[1]{孟子言人但破碎瓦，畫地則復墁滅之，此無用之爲也，然而其意反欲求食，則可食乎？}

曰：「否。」\jiazhu[1]{彭更曰：不食也。}

曰：「然則子非食志也，食功也。」\jiazhu[1]{孟子曰：如是，子果食功也。}\par \jiazhu[1]{章指言：百工食力，以禄養賢，脩仁尚義，國之所尊，移風易俗，其功可珍。雖食諸侯，不爲素餐。}

\section{萬章問宋行王政}

萬章問曰：「宋，小國也。今將行王政，齊、楚惡而伐之，則如之何？」\jiazhu[1]{問宋當如齊、楚何也。}

孟子曰：「湯居亳，與葛爲鄰。葛伯放而不祀，湯使人問之曰：『何爲不祀？』曰：『無以供犧牲也。』湯使遺之牛羊。葛伯食之，又不以祀。」\jiazhu[1]{葛，夏諸侯，嬴姓之國。放縱無道，不祀先祖。}

湯又使人問之曰：「何爲不祀？」曰：「無以供粢盛也。」湯使亳衆往爲之耕，老弱饋食。葛伯率其民，要其有酒食黍稻者奪之，不授者殺之。有童子以黍肉餉，殺而奪之。《書》曰：「葛伯仇餉。」此之謂也。」\jiazhu[1]{童子未成人，殺之尤無狀。書，尚書逸篇也。仇，怨也。言湯所以伐殺葛伯，怨其害此餉也。}

「爲其殺是童子而征之，四海之內皆曰：『非富天下也，爲匹夫匹婦復讎也。』」\jiazhu[1]{四海之民皆曰：湯不貪天下富也，爲一夫報仇也。}

「湯始征，自葛載，十一征而無敵於天下。東面而征，西夷怨；南面而征，北狄怨；曰：『奚爲後我？』民之望之，若大旱之望雨也。歸市者弗止，芸者不變，誅其君，弔其民，如時雨降。民大悅。《書》曰：『徯我后，后來其無罰。』」\jiazhu[1]{載，始也。言湯初征從葛始也。十一征而服天下。一說言當作再字。再十一者，湯再出征十一國。再十一凡征二十二國也。書，逸篇也。民曰待我君來，我則無罰矣。歸市不止，不以有軍來征，故市者止不行也。不使芸者變休也。}

「有攸不惟臣，東征，綏厥士女。匪厥玄黃，紹我周王見休，惟臣附于大邑周。其君子實玄黃于匪以迎其君子，其小人簞食壺漿以迎其小人。救民於水火之中，取其殘而已矣。」\jiazhu[1]{從「有攸」以下，道周武王伐紂時也，皆尚書逸篇之文。攸，所也。言武王東征安天下士女，小人各有所執往，無不惟念執臣子之節。匪厥玄黃，謂諸侯執玄三纁二之帛，願見周王，望見休善，使我得附就大邑周家也。其君子小人各有所執，以迎其類也。言武王之師救殷民於水火之中，討其殘賊也。}

「《大誓》曰：『我武惟揚，侵于之疆，則取于殘，殺伐用張，于湯有光。』」\jiazhu[1]{大誓，古尚書百二十篇之時大誓也。我武王用武之時，惟鷹揚也。侵于之疆，侵紂之疆界。則取于殘賊者，以張伐殺之功也。民有簞食壺漿之歡，比於湯伐桀爲有光寵，美武王德優前代也。今之尚書大誓篇後得以充學，故不與古大誓同。諸傳記引大誓，皆古大誓也。}

「不行王政云爾；苟行王政，四海之內皆舉首而望之，欲以爲君。齊、楚雖大，何畏焉？」\jiazhu[1]{萬章憂宋迫於齊、楚，不得行政，故孟子爲陳殷湯、周武之事以喻之。誠能行之，天下思以爲君，何畏齊、楚焉？}\par \jiazhu[1]{章指言：脩德無小，暴慢無強。是故夏商之末，民思湯武，雖欲不王，末由也已。}

\section{孟子謂戴不勝}

孟子謂戴不勝曰：「子欲子之王之善與？我明告子。」\jiazhu[1]{不勝，宋臣。}

「有楚大夫於此，欲其子之齊語也，則使齊人傅諸？使楚人傅諸？」\jiazhu[1]{孟子假喻：有楚大夫在此，欲變其子使學齊言，當使齊人傳之邪？使楚人自傳相之邪？}

曰：「使齊人傳之。」\jiazhu[1]{不勝曰：使齊人。}

曰：「一齊人傅之，衆楚人咻之，雖日撻而求其齊也，不可得矣。引而置之莊嶽之間數年，雖日撻而求其楚，亦不可得矣。」\jiazhu[1]{言使一齊人傳相，楚衆人咻之。咻之者，嚾也。如此雖日撻之，欲使齊言不可得矣。言寡不勝衆也。莊、嶽，齊街里名也。多人處之數年，而自齊也。}

「子謂薛居州，善士也，使之居於王所。在於王所者，長幼卑尊皆薛居州也，王誰與爲不善？」\jiazhu[1]{孟子曰：不勝常言居州宋之善士也。欲使居於王所。如使在王所者，小大皆如居州，則王誰與爲不善也？}

「在王所者，長幼卑尊皆非薛居州也，王誰與爲善？一薛居州，獨如宋王何？」\jiazhu[1]{如使在王左右者皆非居州之儔，王當誰與爲善乎？一薛居州獨如宋王何而能化之也？周之末世，列國皆僭號自稱王，故曰宋王也。}\par \jiazhu[1]{章指言：自非聖人，在所變化。故諺曰：白沙在涅，不染自黑；蓬生麻中，不扶自直。言輔之者衆也。}

\section{公孫丑問不見諸侯}

公孫丑問曰：「不見諸侯，何義？」\jiazhu[1]{丑怪孟子不肯每輒應諸侯之聘，不見之，於義謂何也？}

孟子曰：「古者不爲臣不見。」\jiazhu[1]{古者不爲臣不肯見，不義而富且貴者也。}

「段干木踰垣而辟之，泄柳閉門而不內，是皆已甚。迫，斯可以見矣。」\jiazhu[1]{孟子言魏文侯、魯繆公有好善之心，而此二人距之大甚。迫窄則可以見之。}

「陽貨欲見孔子而惡無禮。大夫有賜於士，不得受於其家，則往拜其門。」\jiazhu[1]{陽貨，魯大夫也。孔子，士也。}

「陽貨矙孔子之亡也，而饋孔子蒸豚。孔子亦矙其亡也，而往拜之。當是時，陽貨先，豈得不見？」\jiazhu[1]{矙，視也。陽貨視孔子亡而饋之者，欲使孔子來答，恐其便答拜使人也。孔子矙其亡者，心不欲見陽貨也。論語曰饋孔子豚，孟子曰蒸豚。豚，非大牲，故用熟饋也。是時陽貨先加禮，豈得不往拜見之哉？}

「曾子曰：『脅肩諂笑，病于夏畦。』」\jiazhu[1]{脅肩，竦體也。諂笑，強笑也。病，極也。言其意苦勞極，甚於仲夏之月治畦灌園之勤也。}

「子路曰：『未同而言，觀其色赧赧然，非由之所知也。』」\jiazhu[1]{未同，志未合也。不可與言而與之言，謂之失言也。觀其色赧報然，面赤心不正貌也。由，子路名。子路剛直，故曰非由所知也。}

「由是觀之，則君子之所養，可知已矣。」\jiazhu[1]{孟子言由是觀君子子路之言，以觀君子之所養志，可知矣。謂君子養正氣，不以入邪也。}\par \jiazhu[1]{章指言：道異不謀，迫斯強之。段、泄已甚，矙亡得宜。正已直行，不納於邪。赧然不接，傷若夏畦也。}

\section{戴盈之請輕稅}

戴盈之曰：「什一，去關市之征，今兹未能。請輕之，以待來年，然後已，何如？」\jiazhu[1]{戴盈之，宋大夫。問孟子欲使君去關市征稅，復古行什一之賦。今年未能盡去，且使輕之，待來年然後復古，何如？}

孟子曰：「今有人日攘其鄰之雞者，或告之曰：『是非君子之道。』曰：『請損之，月攘一雞，以待來年，然後已。』如知其非義，斯速已矣，何待來年？」\jiazhu[1]{攘，取也，取自來之物也。孟子以此爲喻，知攘之惡當即止，何可損少，月取一雞，待來年乃止乎？謂盈之之言若此類者也。}\par \jiazhu[1]{章指言：從善改非，坐而待旦。知而爲之，罪重於故。譬猶攘雞，多少同盜，變惡自新，速然後可也。}

\section{公都子問好辯}

公都子曰：「外人皆稱夫子好辯，敢問何也？」\jiazhu[1]{公都子，孟子弟子也。外人，他人論議者。好辯，言子好與揚、墨之徒辯爭。}

孟子曰：「予豈好辯哉？予不得已也。」\jiazhu[1]{曰：我不得已耳，欲救正道，懼爲邪說所亂，故辯之也。}

「天下之生久矣，一治一亂。當堯之時，水逆行，氾濫於中國，蛇龍居之，民無所定；下者爲巢，上者爲營窟。」\jiazhu[1]{天下之生，生民以來也。迭有亂治，非一世。水生蛇龍，水盛則蛇龍居民之地也。民患水避之，故無定居。埤下者於樹上爲巢，猶鳥之巢也。上者，高原之上也，鑿岸而營度之，以爲窟穴而處之。}

「《書》曰：『洚水警余。』洚水者，洪水也。」\jiazhu[1]{尚書逸篇也。水逆行，洚洞無涯，故曰洚水。洪，大也。}

「使禹治之。禹掘地而注之海，驅龍蛇而放之菹。水由地中行，江、淮、河、漢是也。險阻既遠，鳥獸之害人者消，然後人得平土而居之。」\jiazhu[1]{堯使禹治洪水通九州，故曰掘地而注之海也。菹，澤生草者也，今青州謂澤有草者爲菹。水流行於地而去也。民人下高就平土，故遠險阻也。水去故鳥獸害人者消盡也。}

「堯、舜既沒，聖人之道衰，暴君代作。壞宮室以爲汙池，民無所安息；棄田以爲園囿，使民不得衣食。邪說暴行又作，園囿、汙池、沛澤多而禽獸至。」\jiazhu[1]{暴，亂也。亂君更興，殘壞民室屋，以其處爲汙池，棄五穀之田以爲園囿，長逸遊而棄本業，使民不得衣食，有飢寒竝至之厄。其小人則放辟邪侈，故作邪僞之說，爲姦寇之行。沛，草木之所生也。澤，水也。至，衆也。田疇不墾，故禽獸衆多，謂羿、桀之時也。}

「及紂之身，天下又大亂。周公相武王誅紂，伐奄三年討其君，驅飛廉於海隅而戮之，滅國者五十，驅虎、豹、犀、象而遠之，天下大悅。」\jiazhu[1]{奄，東方無道國。武王伐紂至於孟津還歸，二年復伐，前後三年也。飛廉，紂諛臣。驅之海隅而戮之，猶舜放四罪也。滅與紂共爲亂政者五十國也。奄，大國，故特伐之。尚書多方曰：「王來自奄。」}

「《書》曰：『丕顯哉，文王謨！丕承哉，武王烈！佑啓我後人，咸以正無缺。』」\jiazhu[1]{書，尚書逸篇也。丕，大；顯，明；承，纘；烈，光也。言文王大顯明王道，武王大纘承天光烈，佑開後人，謂成、康皆行正道無虧缺也。此周公輔相以撥亂之功也。}

「世衰道微，邪說暴行有作，臣弒其君者有之，子弒其父者有之。孔子懼，作《春秋》。《春秋》，天子之事也。是故孔子曰：『知我者其惟《春秋》乎！罪我者其惟《春秋》乎！』」\jiazhu[1]{世衰道微，周衰之時也。孔子懼王道遂滅，故作《春秋》，因魯史記，設素王之法，謂天子之事也。知我者，謂我正王綱也；罪我者，謂時人見彈貶者。言孔子以《春秋》撥亂也。}

「聖王不作，諸侯放恣，處士橫議，楊朱、墨翟之言盈天下。天下之言不歸楊，則歸墨。楊氏爲我，是無君也；墨氏兼愛，是無父也。無父無君，是禽獸也。」\jiazhu[1]{言孔子之後，聖人之道不興，戰國縱橫，布衣處士游說以干諸侯，若楊、墨之徒，無尊異君父之義，而以橫議於世也。}

「公明儀曰：『庖有肥肉，廏有肥馬，民有飢色，野有餓莩，此率禽獸而食人也。』」\jiazhu[1]{公明儀，魯賢人。言人君但崇庖廚，養犬馬，不恤民，是爲率禽獸而食人也。}

「楊、墨之道不息，孔子之道不著，是邪說誣民，充塞仁義也。仁義充塞，則率獸食人，人將相食。」\jiazhu[1]{言仁義塞則邪說行，獸食人則人相食，此亂之甚也。}

「吾爲此懼，閑先聖之道，距楊、墨，放淫辭，邪說者不得作。」\jiazhu[1]{閑，習也；淫，放也。孟子言我懼聖人之道不著，爲邪說所乘，故習聖人之道以距之。}

「作於其心，害於其事；作於其事，害於其政。聖人復起，不易吾言矣。」\jiazhu[1]{說與上篇同。}

「昔者禹抑洪水而天下平，周公兼夷狄、驅猛獸而百姓寧，孔子成《春秋》而亂臣賊子懼。」\jiazhu[1]{抑，治也。周公兼懷夷狄之人，驅害人之猛獸也。言亂臣賊子懼《春秋》之貶責也。}

「《詩》云：『戎狄是膺，荊舒是懲，則莫我敢承。』」\jiazhu[1]{此詩巳見上篇說。}

「無父無君，是周公所膺也。」\jiazhu[1]{是周公所欲伐擊也。}

「我亦欲正人心，息邪說，距詖行，放淫辭，以承三聖者。豈好辯哉？予不得已也。」\jiazhu[1]{孟子言我亦欲正人心，距險詖之行，以奉禹、周公、孔子也。不得已而與人辯耳，豈好之哉？}

「能言距楊、墨者，聖人之徒也。」\jiazhu[1]{孟子自謂能距楊、墨也。徒，黨也。可以繼聖人之道，謂名世者也。}\par \jiazhu[1]{章指言：夫憂世撥亂，勤以濟之，義以匡之，是故禹、稷駢躓，周公仰思，仲尼皇皇，墨突不及汙。聖賢若是，豈得不辯也？}

\section{匡章論陳仲子}

匡章曰：「陳仲子豈不誠廉士哉？居於陵，三日不食，耳無聞，目無見也。井上有李，螬食實者過半矣，匍匐往，將食之，三咽，然後耳有聞，目有見。」\jiazhu[1]{匡章，齊人也。陳仲子，齊一介之士，窮不苟求者，是以絕糧而餒也。螬，蟲也。李實有蟲食之過半，言仲子目不能擇也。}

孟子曰：「於齊國之士，吾必以仲子爲巨擘焉。雖然，仲子惡能廉？充仲子之操，則蚓而後可者也。夫蚓，上食槁壤，下飲黃泉。」\jiazhu[1]{巨擘，大指也。比於齊國之士，吾必以仲子爲指中大者耳，非大器也。蚓，丘蚓之蟲也。充滿其操行，似蚓而可行者也。蚓食土飲泉，極廉矣，然無心無識。仲子不知仁義，苟守一介，亦猶蚓也。}

「仲子所居之室，伯夷之所築與？抑亦盜跖之所築與？所食之粟，伯夷之所樹與？抑亦盜跖之所樹與？是未可知也。」\jiazhu[1]{孟子問匡章：仲子豈能必使伯夷之徒築室樹粟，乃居食之邪？抑亦得盜跖之徒使作也？是殆未可知也。}

曰：「是何傷哉？彼身織屨，妻辟纑，以易之也。」\jiazhu[1]{匡章曰：惡人作之何傷哉？彼仲子身自織屨，妻緝纑以易食宅耳。緝績其麻曰辟，練麻曰纑。}

曰：「仲子，齊之世家也。兄戴，蓋禄萬鍾。以兄之禄爲不義之禄而不食也，以兄之室爲不義之室而不居也，避兄離母，處於於陵。」\jiazhu[1]{孟子言仲子齊之世卿大夫之家。兄名戴，爲齊卿，食采於蓋，禄萬鍾。仲子以爲事非其君，行非其道，以居富貴，故不義之，竄於於陵。}

「他日歸，則有饋其兄生鵝者。己頻顣曰：『惡用是鶃鶃者爲哉？』」\jiazhu[1]{他日，異日也。歸省其母，見兄受人之鵝而非之。己，仲子也。頻顣，不悅曰：安用是鶃鶃者爲乎？鶃鶃，鵝鳴之聲。}

「他日，其母殺是鵝也，與之食之。其兄自外至，曰：『是鶃鶃之肉也。』出而哇之。」\jiazhu[1]{異日母食以鵝，不知是前所頻顣者也。兄疾之，告曰是鶃鶃之肉也。仲子出門而哇吐之。}

「以母則不食，以妻則食之；以兄之室則弗居，以於陵則居之。是尚爲能充其類也乎？若仲子者，蚓而後充其操者也。」\jiazhu[1]{孟子非其不食於母而食妻所作屨纑易食也，不居兄室而居於於陵人所築室，是尚能充人類乎？如蚓之性，然後可以充其操也。}\par \jiazhu[1]{章指言：聖人之道，親親尚和。志士之操，耿介特立。可以激濁，不可常法。是以孟子喻以丘蚓，比諸巨擘也。}
